(1) 半径 \(r\) 内质量 \(M(r)=4\pi\rho_0\int_0^r(s^2-s^4/R^2)ds=4\pi\rho_0(r^3/3-r^5/(5R^2))\)。壳层定理给 \(g_r=-GM(r)/r^2=-4\pi G\rho_0(r/3-r^3/(5R^2))\)，有向 \(x\) 的加速度相同取 \(r\mapsto x\)。定义 \(\omega_0^2=4\pi G\rho_0/3\)，球心附近受力为 \(-m\omega_0^2x+O(x^3)\)。

(2) 令 \(\alpha=4\pi G\rho_0/(5R^2)\)，精确的内部运动方程是 \(\ddot x+\omega_0^2x-\alpha x^3=0\)。单位质量有效势（以球心零点）为 \(V(x)=\omega_0^2x^2/2-\alpha x^4/4\)，且 \(V'(x)=\omega_0^2x-\alpha x^3\)。它的力学能 \(\dot x^2/2+V(x)\) 在隧道内守恒；非线性项符号为负，表明往大振幅方向恢复力相对减弱。

(3) 采用 Lindstedt 方法设 \(x=A\cos(\omega t)+O(A^3)\)，代入 \(\ddot x+\omega_0^2x-\alpha x^3=0\)；\(\cos^3 u=(3\cos u+\cos3u)/4\)，消去基频共振项得到 \(\omega^2=\omega_0^2-3\alpha A^2/4+O(A^4)\)，即 \(\omega=\omega_0[1-3\alpha A^2/(8\omega_0^2)+O(A^4/R^4)]=\omega_0[1-9A^2/(40R^2)+\cdots]\)。因此周期 \(T=(2\pi/\omega_0)[1+9A^2/(40R^2)+\cdots]\)。

(4) 总质量 \(M(R)=8\pi\rho_0R^3/15\)，故边界势 \(\Phi(R)=-8\pi G\rho_0R^2/15\)。利用 \(\Phi(R)-\Phi(0)=\int_0^R GM(r)/r^2dr=7\pi G\rho_0R^2/15\)，得 \(\Phi(0)=-\pi G\rho_0R^2\)。以无穷远势零为准，球心逃逸条件 \(v_{\mathrm{esc}}^2/2+\Phi(0)=0\)，于是 \(v_{\mathrm{esc}}=R\sqrt{2\pi G\rho_0}\)。注意不能以均匀球体的势替代当前非均匀密度。
